\documentclass[11pt,a4paper]{article}
\usepackage[utf8]{inputenc}\usepackage[T1]{fontenc}\usepackage[english]{babel}
\usepackage{mathptmx}
\usepackage[top=2.2cm,bottom=2.2cm,left=2.5cm,right=2.5cm]{geometry}
\usepackage{setspace}\onehalfspacing
\usepackage{parskip}\usepackage{microtype}
\usepackage{graphicx}\usepackage{float}\usepackage{booktabs}
\usepackage{tabularx}\usepackage{array}\usepackage{amssymb}\usepackage{amsmath}
\usepackage{xcolor}
\usepackage[hidelinks,colorlinks=true,linkcolor=blue!60!black,citecolor=blue!60!black]{hyperref}
\usepackage[numbers,sort&compress]{natbib}
\pagestyle{plain}

\title{\textbf{Can Language Models Simulate Psychometric Respondents?\\[4pt]
A Cross-Vendor Sensitivity Study with Baseline Controls\\and a Replicated Zero-Prompt Framing Manipulation\\for a Polish Mentalization Battery}}
\author{Micha{\l} Wiencek\\\small Independent researcher; MA student, Institute of Psychology, UKEN Krak\'ow}
\date{}

\begin{document}
\maketitle

\begin{abstract}
\noindent Before collecting human data for a new psychometric battery, I test whether the instruments \emph{react} to personality variation at all. Thirty fictional Polish-language biographies with predetermined psychological profiles were processed by five large language models from four vendors --- Claude Sonnet 4.6 and Opus 4.6 (Anthropic), GPT-5.4-mini (OpenAI), Grok-4-20-reasoning (xAI), and Gemini 3 Flash (Google). Each persona was attempted twice per model (300 attempts; 285 valid persona runs after exclusion of partial outputs and content-filtered responses), and 7 human respondents completed the same battery (292 observations including human pilot). This is a \textbf{sensitivity study}, not a validation.

Four published self-report instruments showed consistent profile recovery: attachment $\kappa = .33$--$.87$ across models; rank--score correlations median $r = .71$--$.82$; literature consistency 8--10/10. Cross-vendor questionnaire agreement was high (median $r = .78$--.96), but TCTM-22 cross-vendor agreement was low (ICC $= .25$). Error profiles were model-dependent: Claude Sonnet produced almost exclusively DOS (97\%); Gemini produced the descriptively closest-to-human-pilot error profile.

Two baseline experiments decompose the observed variance into three components. In the \textit{baseline condition} ($N = 46$ valid runs, 7--10 per model), models completed the battery with an ``answer as yourself'' instruction, revealing stable ``default response profiles'': GPT is anxious-preoccupied in 100\% of runs, Gemini hits near-maximum on all self-report scales, and different models from the same vendor (Sonnet vs.\ Opus) differ in attachment style. The \textit{zero-prompt condition} was extended to a replication of $N = 5$--$8$ runs per model: 31 valid TCTM-22 payloads and 30 parseable self-report payloads (one GPT zero-prompt output had a complete TCTM-22 array but no parseable self-report sections). With multiple observations the ``answer as yourself'' instruction provides replicated evidence of systematic shifts in self-report profiles: Sonnet's avoidance dropped from baseline $M = 3.88$ to zero-prompt $M = 2.98$ ($N = 6$, $\Delta = -0.90$, closely replicating the original single-shot estimate of $-0.99$); Gemini's emotional stability dropped from baseline $M = 7.0$ to zero-prompt $M = 6.0$ ($N = 8$, $\Delta = -1.00$, exact replication); and four of five models reported \emph{higher} attachment anxiety without the instruction (Sonnet $+1.02$, Opus $+0.55$, Grok $+0.75$, Gemini $+0.90$). Author-keyed TCTM-22 accuracy was preserved across framing conditions for four of five models; GPT dropped from baseline $78.6\%$ to zero-prompt $70.5\%$ ($-8.1$ pp), the only model showing a notable accuracy shift. Persona manipulation increased response variance 2--51$\times$ over baseline stochasticity (model-level medians 3.5--8.9$\times$ across 9 questionnaire dimensions where baseline $SD > 0$), confirming that persona-driven variance is not an artifact of inherent model noise.

Core materials (biographies, scored responses, analysis scripts) are released as an open dataset.

\vspace{4pt}\noindent\textbf{Keywords:} synthetic respondents, instrument sensitivity, mentalization, attachment, cross-vendor LLM comparison, baseline control, LLM default response profile
\end{abstract}

% ═══════════════════════════════════
\section{Introduction}

Developing new psychometric instruments requires extensive piloting. A growing line of research suggests that large language models (LLMs) can serve as \emph{simulated respondents}, providing rapid probes of instrument behavior before human data collection \cite{argyle2023,dillion2023}. The core question is not whether LLMs ``have'' personality, but whether instruments \emph{react systematically} to known personality variation embedded in biographical prompts. If they do, the instruments are at minimum sensitive to the relevant signal --- a necessary (though far from sufficient) precondition for construct validity \cite{cronbach1955}. If they do not, there is a problem worth catching before human recruitment.

This framing is important: the present study is a \textbf{sensitivity test}, not a validation study in the classical sense. I ask: ``Given 30 biographies designed to span the full range of attachment styles, mentalization levels, and personality traits, can the instruments recover the embedded profiles from LLM responses?'' A positive answer means the instruments are responsive; it does not prove they measure the intended constructs in human respondents.

Most existing studies use a single model family \cite{argyle2023,dillion2023}; by contrast, the present study employs \textbf{five models from four vendors} (Anthropic, OpenAI, xAI, Google), enabling assessment of cross-vendor consistency --- whether sensitivity findings generalize across architectures or are model-specific artifacts. Crucially, two baseline experiments are included: a \textit{baseline condition} (models complete the battery ``as themselves,'' repeated 7--10 times each) and a \textit{zero-prompt condition} (models receive the battery with no system instruction at all). This design allows decomposition of the observed persona-driven variance into three components: inherent model stochasticity, the effect of instructional framing, and the genuine persona manipulation effect.

I apply this approach to four published Polish-language instruments --- DBZ-R (attachment), MentS-PL (mentalization), KPP (need for cognition), TIPI-PL (Big Five) --- and one original mentalization vignette test (TCTM-22). Thirty fictional personas were constructed, processed through all five models with test--retest, and scored against ground-truth profiles. Seven human respondents completed the same battery for exploratory comparison.

% ═══════════════════════════════════
\section{Method}

\textbf{Personas.} 30 biographies in Polish (1,554--2,914 file-total words including YAML metadata; 1,489--2,861 narrative-body words after stripping YAML), written in second person, encoding predetermined profiles across 12 dimensions (11 continuous + 1 categorical): attachment anxiety and avoidance (DBZ-R), mentalization self/other/motivation (MentS-PL), need for cognition (KPP), Big Five E/A/C/ES/O (TIPI-PL), and attachment style (4-category classification). Biographies were authored by the first author using Claude Opus 4.6 as a writing assistant; the author specified profiles, narrative constraints, and stylistic requirements, reviewed and edited all drafts. Six biographies (\texttt{adrian, ewa, kuba, magda, pawel, weronika}) were revised after pilot review and are tagged \texttt{persona\_version: v2} in the released YAML; the remaining 24 are at \texttt{v1}. The version that was actually sent to all five models is the version referenced throughout this paper and recorded in \texttt{all\_data.csv}. No biography was designed to contain trait labels; the first author manually coded biographies into two groups based on whether they contained any direct trait-language references (e.g., explicit mentions of attachment style, introversion, openness): nine biographies were flagged as ``explicit'' (\texttt{ania, anna-sim, filip, jakub, kamil, marek, natalia, ola, tomek}) and 21 as ``narrative''. A separate regex-based audit using a broader keyword set yields a different split (see Trait Transparency in Results); both codings are released in the supplementary archive.

\textbf{Instruments.} DBZ-R (Polish adaptation of ECR-R; 36 items, Likert 1--7; Lubiewska et al., 2016), MentS-PL (28 items; Ja\'nczak, 2021), KPP (36 items; Matusz et al., 2011), TIPI-PL (10 items; Sorokowska et al., 2014). These four instruments use published Polish norms. TCTM-22 (22 original vignettes with MASC error taxonomy \cite{dziobek2006}) is an original instrument without published norms; it is scored against author-defined answer keys. The 22 vignettes cover five subscales (4--5 items each): \emph{Rozpoznawanie podtekstu} (subtext recognition; 4 items: w01, s07, s08, s10), \emph{Rozpoznawanie manipulacji} (manipulation recognition; 4 items: w08, c07, c10, w11), \emph{Wykrywanie sygna\l\'ow wygasania wi\k{e}zi} (waning-bond detection; 4 items: w13, w14, e08, w15), \emph{Dynamika d\k{a}\.zenie--wycofanie} (pursuit--withdrawal dynamics; 5 items: w19, pw07, w22, pw09, pw11), and \emph{Ocena autentyczno\'sci naprawy relacyjnej} (authenticity of relational repair; 5 items: w25, r08, w28, r09, r10).

\textbf{Ground truth and STYLE\_MAP.} Each persona's expected attachment style is taken from the top-level \texttt{expected\_style} field in \texttt{analysis\_summary.jsonl}. Four personas (Ewa, Kamil, Marek, Radek) are coded as \texttt{disorganized} in the YAML \texttt{expected\_profile} block of their biography file but as \texttt{fearful\_avoidant} in the top-level \texttt{expected\_style} field used by the scoring pipeline; this is intentional and corresponds to mapping disorganized attachment to fearful-avoidant under the four-category Bartholomew taxonomy that all five LLMs use in their style outputs. Computing $\kappa$ under a strict five-category policy (treating \texttt{disorganized} as a separate fifth class that no model can produce) would lower all attachment $\kappa$ values by $\sim 10$ percentage points; the supplementary archive includes both calculations.

\textbf{Models.} Five models from four vendors (Table~\ref{tab:models}). Each persona was processed once per model (run~1), then re-run under identical conditions (run~2, retest). All models used default sampling parameters. Total: 285 valid persona runs (300 attempted; missing retests resulted from API content filters, partial generations, or incomplete TCTM responses --- see Table~\ref{tab:models} footnote). Plus 5 additional Gemini reruns (run~3) for personas \texttt{ania, anna-sim, kasia, marek, michal-sim} where the first two attempts returned partial or degenerate output; these run-3 records are included in the released open dataset for full transparency but are excluded from all run~1 / run~2 paired analyses, all cross-model first-run statistics, and all main-text tables and figures, to avoid asymmetric over-representation of Gemini. GPT-5.4-mini returned incomplete TCTM responses (20--21 instead of 22 items) in 20\% of attempts; affected runs were excluded from per-model TCTM tables (yielding $N = 27$ for run~1, $N = 21$ for run~2). For other instruments, missing items were scored according to the published convention (DBZ-R / MentS-PL / KPP / TIPI-PL: missing-as-NA, with row-wise pairwise deletion in correlation analyses).

\begin{table}[htbp]\centering\footnotesize
\caption{Models used. Run counts reflect valid completions after exclusion of partial / content-filtered attempts. The 5 additional Gemini run-3 records are partial-fix retries for personas with earlier degenerate outputs and are reported transparently rather than discarded.}
\label{tab:models}
\begin{tabular}{llll}
\toprule
\textbf{Model} & \textbf{Vendor} & \textbf{Platform} & \textbf{Runs (run~1 + run~2)} \\
\midrule
Claude Sonnet 4.6 & Anthropic & AWS Bedrock & $30 + 30$ \\
Claude Opus 4.6 & Anthropic & AWS Bedrock & $30 + 25^{\ddagger}$ \\
GPT-5.4-mini & OpenAI & Azure OpenAI & $27^{\dagger} + 21^{\dagger}$ \\
Grok-4-20-reasoning & xAI & Azure AI Foundry & $30 + 29^{\ddagger}$ \\
Gemini 3 Flash & Google & Gemini API & $30 + 28^{\ddagger}$ ($+5$ run-3) \\
\midrule
Human respondents & --- & --- & 7 \\
\bottomrule
\end{tabular}
\begin{flushleft}
\footnotesize $^{\dagger}$Azure content filter or partial TCTM response. $^{\ddagger}$One or more personas failed run-2 due to transient API issues (timeout, refusal, or partial JSON). For Opus run-2 the missing personas are \texttt{anna-sim, jakub, kasia, michal-sim, piotr}; for Grok run-2 it is \texttt{piotr}; for Gemini run-2 it is \texttt{ewa, sara}. None were retried under the v4 protocol; the missing-data pattern is disclosed here for full transparency.
\end{flushleft}
\end{table}

\textbf{Human pilot.} Seven respondents (ages 22--60, 4 male, 3 female; 2 psychology students, 5 without psychology background) completed the same battery online. The target population for eventual validation is young adults (ages 18--35) in romantic or close relationships; one respondent (R07, age 60, female) was notably older than this range and was retained for exploratory purposes.

\textbf{Ethics.} The human pilot used informed consent disclosing study purpose, voluntariness (right to withdraw at any point without consequence), GDPR-compliant data handling (encrypted storage, five-year retention, no cross-linking with personal identifiers), and links to mental-health support resources (Polish national crisis-line numbers). All seven respondents consented before completing the battery. The pilot is non-interventional, uses no clinical population, and was treated as exempt from formal Institutional Review Board approval per UKEN Krak\'ow institutional guidelines for pre-pilot psychometric data collection; this exemption status reflects the author's understanding of those guidelines and is disclosed here for full transparency.

\textbf{Baseline condition.} Each model completed the battery 7--10 times (46 valid runs total; one Gemini run excluded for incomplete TCTM) with the system instruction: \textit{``You are an AI language model. Answer the questions honestly, as you --- an AI --- would answer if you had to choose. Do not play any character. Answer as yourself.''} (translated from Polish). This condition measures inherent model stochasticity and reveals each model's ``default response profile.''

\textbf{Zero-prompt condition.} Each model completed the battery $N = 5$--$8$ times without any system instruction --- only the test content was provided as a user message: 31 valid TCTM-22 payloads (Sonnet 6, Opus 5, GPT 6, Grok 6, Gemini 8) and 30 parseable self-report payloads (one GPT zero-prompt output had a complete TCTM-22 array but no parseable self-report sections). The same TCTM-22 instrument used in the persona condition was administered, alongside the four self-report instruments, in a single user message. Without instructional context, models produced widely varying JSON formats (e.g., Sonnet alternated between \texttt{TCTM22}, \texttt{czesc\_1\_winiety}, and \texttt{part1\_vignettes} as top-level keys across runs); responses were post-processed with a universal scorer that filters to the 22 standard TCTM items and applies the validated \texttt{score\_dbzr/ments/kpp/tipi} functions used elsewhere in this study. This condition measures the model's ``raw'' response to a psychometric task, before any instructional framing.

\textbf{Metrics.} (1) Attachment style match ($\kappa$, Cram\'er's $V$), (2) rank--score correlations (Pearson, Spearman), (3) literature consistency (10 construct pairs), (4) test--retest reliability (Pearson $r$, ICC), (5) inter-model consistency (ICC, Fleiss' $\kappa$), (6) TCTM-22 error profiles (DOS/NAD/BK), (7) per-item cross-vendor agreement (Cochran's $Q$).

% ═══════════════════════════════════
\section{Results}

All cross-model tables report first-run data ($N = 30$ per model except where the per-model footnote on Tables~\ref{tab:tctm} and \ref{tab:noprompt-tctm} indicates a smaller valid-output count; specifically, GPT-5.4-mini run~1 has $N = 27$ for analyses requiring a complete 22-item TCTM-22 response). Retest data is used for reliability analyses and, where noted, to confirm systematic patterns.

\subsection{TCTM-22 Score Distribution}

\begin{table}[htbp]\centering\footnotesize
\caption{TCTM-22 by group (first run).}
\label{tab:tctm}
\begin{tabular}{lcccccccc}
\toprule
\textbf{Group} & $N$ & $M$ & $SD$ & Range & DOS\% & NAD\% & BK\% \\
\midrule
Human       & 7  & 14.3 & 1.4 & 12--16 & 32 & 36 & 32 \\
GPT-5.4m    & 27$^\dagger$ & 17.0 & 1.7 & 13--19 & 70 & 6 & 24 \\
Grok-4-20   & 30 & 18.4 & 4.2 & 6--22  & 64 & 16 & 21 \\
Opus        & 30 & 19.0 & 2.1 & 8--20  & 86 & 8 & 7 \\
Sonnet      & 30 & 19.4 & 0.7 & 17--20 & 97 & 0 & 3 \\
Gemini 3F   & 30 & 19.6 & 2.1 & 9--21  & 23 & 41 & 37 \\
\bottomrule
\end{tabular}
\begin{flushleft}
\footnotesize $^\dagger$GPT returned no valid 22-item response for 3 of 30 personas (Pawe{\l}, Micha{\l}-K, Bartek).
\end{flushleft}
\end{table}

\begin{figure}[htbp]\centering
\includegraphics[width=0.85\textwidth]{synthetic/charts/v5_tctm_boxplot.png}
\caption{TCTM-22 distribution by group: Human + 5 LLM models. Diamond = mean.}
\end{figure}

A clear gradient emerged: Human ($M = 14.3$) $<$ GPT ($17.0$) $<$ Grok ($18.4$) $<$ Opus ($19.0$) $\approx$ Sonnet ($19.4$) $\approx$ Gemini ($19.6$).

\begin{figure}[htbp]\centering
\includegraphics[width=0.85\textwidth]{synthetic/charts/v5_error_profile.png}
\caption{Error type distribution: Human + 5 LLM models. Humans show a roughly balanced profile (DOS 32\%, NAD 36\%, BK 32\%); Claude Sonnet is almost exclusively DOS (97\%). Among LLMs, Gemini's profile (DOS 23\%, NAD 41\%, BK 37\%) is closest to the human distribution.}
\end{figure}

Error profiles were model-dependent. Claude Sonnet produced almost exclusively DOS errors (97\%); Gemini produced the closest match to the human pilot profile (DOS 23\%, NAD 41\%, BK 37\% vs.\ human 32/36/32\%). In the human pilot, NAD errors were more frequent than in all LLMs except Gemini, which exceeded the human NAD rate (41\% vs.\ 36\%).

\subsection{Attachment Classification}

\begin{table}[htbp]\centering\footnotesize
\caption{Attachment style classification.}
\begin{tabular}{lccc}
\toprule
\textbf{Model} & \textbf{Accuracy} & $\kappa$ & $V$ \\
\midrule
Opus        & 27/30 = 90\% & .866 & .905 \\
Grok-4-20   & 27/30 = 90\% & .867 & .890 \\
Gemini 3F   & 25/30 = 83\% & .778 & .801 \\
Sonnet      & 24/30 = 80\% & .735 & .765 \\
GPT-5.4m    & 13/27$^{\dagger}$ = 48\% & .333 & .474 \\
\bottomrule
\end{tabular}
\begin{flushleft}
\footnotesize $^{\dagger}$GPT~run-1 with valid 22-item TCTM only ($N = 27$); the 3 personas with missing TCTM (Pawe{\l}, Micha{\l}-K, Bartek) are excluded from $\kappa$ computation rather than scored as misclassifications.
\end{flushleft}
\end{table}

\begin{figure}[htbp]\centering
\includegraphics[width=0.55\textwidth]{synthetic/charts/v5_attachment_kappa.png}
\caption{Attachment $\kappa$ by model (5 LLMs).}
\end{figure}

\subsection{Rank--Score Correlations}

Median Pearson $r$ across 11 continuous dimensions (first run): Opus $.819$, Grok $.815$, Sonnet $.800$, Gemini $.763$, GPT $.711$. Across all five models, 54/55 dimension-level correlations were positive (expected direction); the single negative case was GPT on MentS-Self ($r = -.06$), where higher expected mentalization-self failed to track higher self-reported scores. Separately, literature consistency (10 published cross-scale correlation signs): Gemini 10/10, Sonnet 10/10, Opus 9/10, GPT 9/10, Grok 8/10.

\subsection{Test--Retest Reliability}

\begin{table}[htbp]\centering\footnotesize
\caption{Test--retest reliability (run~1 vs.\ run~2).}
\label{tab:retest}
\begin{tabular}{lccccc}
\toprule
\textbf{Metric} & \textbf{Sonnet} & \textbf{Opus} & \textbf{GPT} & \textbf{Grok} & \textbf{Gemini} \\
\midrule
$z$-score median $r$ & .981 & .993 & .890 & .950 & .962 \\
TCTM total $r$ & .691 & .992 & .279 & .530 & .905 \\
Style agreement & 29/30 & 25/25 & 18/21 & 29/29 & 25/28 \\
$N$ paired personas & 30 & 25 & 21 & 29 & 28 \\
\bottomrule
\end{tabular}
\end{table}

\begin{figure}[htbp]\centering
\includegraphics[width=0.65\textwidth]{synthetic/charts/v5_retest_5model.png}
\caption{Test--retest reliability across 5 models. Solid bars: TCTM-22 total $r$. Lighter bars: $z$-score median $r$ across 9 questionnaire dimensions. Reliability is computed on \emph{intersect-only paired personas} (complete-pair policy, $N$ varies by model).}
\end{figure}

\noindent Opus was near-deterministic across all tasks. Gemini showed high TCTM reliability ($r = .91$). GPT was unreliable on TCTM ($r = .28$ on the 21-persona intersect of valid run-1/run-2 pairs), with individual personas shifting by up to $\pm 5$ points between runs; its questionnaire data was more stable ($z$-score median $r = .89$). All test--retest correlations are computed under a complete-pair policy: paired correlations use only personas with valid responses in both runs (intersect $N$ varies by model and is shown in the table), rather than imputing missing pairs.

\subsection{Cross-Vendor Consistency}

\begin{figure}[htbp]\centering
\includegraphics[width=0.55\textwidth]{synthetic/charts/v5_intermodel_matrix.png}
\caption{Inter-model $z$-score median $r$ (9 questionnaire dimensions, 5$\times$5 matrix). Diagonal $= 1$.}
\end{figure}

Questionnaire $z$-scores showed high cross-vendor agreement: median inter-model $r$ ranged from $.78$ (GPT$\times$Gemini) to $.96$ (Sonnet$\times$Opus). TCTM-22 cross-vendor agreement was low: ICC(2,1) $= .254$ across 5 models on the $N = 27$ intersect of complete TCTM data. The highest TCTM pair was Gemini$\times$Opus ($r = .94$); the lowest was GPT$\times$Grok ($r = -.05$, effectively zero with a slight negative sign on the intersect).

Fleiss' $\kappa$ for attachment style classification (all 5 models as raters, $N = 27$ personas with complete data across all models): $.619$ (``substantial agreement'' under the widely-used categorical-data benchmarks of \cite{landis1977}; we note that those benchmarks have been criticized as arbitrary, and report the raw $\kappa$ value to permit reader-specific interpretation). Cochran's $Q$ revealed that 14/22 TCTM items showed significant ($p < .05$, uncorrected, $df = 4$) model disagreement. Under Bonferroni correction at $\alpha/22 = .00227$ (chi-square threshold $> 16.64$), 8/22 items remain significant: w01, s07, w11, w15, pw07, w22, w28, r10 (all with $Q > 32$, $p < 10^{-5}$). Under Benjamini--Hochberg FDR control at $q = .05$, 13/22 items are significant. The full per-item $Q$, $p_{\text{unc}}$, Bonferroni-pass and BH-pass values are released in the supplementary archive.

\subsection{Ola: Cross-Vendor Item Agreement}

Persona Ola (``secure but non-mentalizing'') had the lowest cross-model mean TCTM ($11.4/22$: Sonnet 17, Opus 8, GPT 16, Grok 7, Gemini 9). Cross-vendor error analysis revealed \textbf{systematic item-level agreement}: all 5 models produced DOS errors on items s07 and w22. This pattern was confirmed by retest on Opus ($8 \rightarrow 9/22$) and Gemini ($9 \rightarrow 7/22$).

\begin{figure}[htbp]\centering
\includegraphics[width=0.85\textwidth]{synthetic/charts/v5_convergence_scatter.png}
\caption{TCTM $\times$ MentS-Total by model. Convergence varies in both magnitude and sign: Grok $r = +.51$ (highest), Opus $r = +.24$ (weakly positive), Gemini $r = +.12$, Sonnet $r = +.01$ (near-zero), GPT $r = -.29$ (weakly \emph{negative}: in GPT, higher self-reported mentalization is associated with lower TCTM accuracy).}
\end{figure}

\begin{figure}[htbp]\centering
\includegraphics[width=0.95\textwidth]{synthetic/charts/v5_item_heatmap.png}
\caption{Per-item accuracy: Human + 5 LLM models. Green = high, red = low.}
\end{figure}

\subsection{Exploratory Human Comparison ($N = 7$)}

Seven human respondents (ages 22--60) scored $M = 14.3/22$ ($SD = 1.4$), well below all LLM means. Human error profiles were distributed across DOS (32\%), NAD (36\%), and BK (32\%). On most author-keyed TCTM-22 items, LLM scores were descriptively higher than the small human pilot: on 7/22 items (w08, c10, w13, e08, w19, w25, w28) the cross-model LLM mean was $\geq 30$ percentage points above human accuracy (range $+34$ to $+65$ pp). On only 2/22 items did human scores exceed LLM scores by a comparable margin: w22 (human 57\% vs.\ Sonnet/Opus 0\%, GPT 70\%, Grok 47\%, Gemini 13\% --- a Claude-specific failure rather than a general LLM ceiling) and s07 (human 71\% vs.\ Sonnet 0\%, Opus 7\%, GPT 0\%, Grok 30\% --- where Gemini is a notable outlier scoring 87\%, descriptively above human accuracy). These observations are descriptive only --- with $N = 7$ humans, no statistical inference is warranted.

\subsection{Trait Transparency}

Composite fidelity was virtually identical for narrative ($n = 21$, fidelity $= 0.790$) and explicit ($n = 9$, fidelity $= 0.783$; $\Delta = -0.007$). Directional accuracy was +4.3pp higher for explicit personas, but untestable at this sample size. \textbf{Important caveat:} the trait-transparency fidelity values reported here are computed from the original 2-model validation pool (Sonnet 4.6 + Opus 4.6, $n = 60$ persona-model observations), generated 12 April 2026 prior to the 5-model expansion of this study. Recomputing fidelity on the full 5-model pool ($n = 285$) was not part of this submission and may yield different values; the original 2-model statistic is retained here for reproducibility with the released \texttt{validation\_report.json}. The narrative/explicit grouping itself is also sensitive to coding choices --- the EXPLICIT category is defined by a hand-curated list of nine biographies (\texttt{ania, anna-sim, filip, jakub, kamil, marek, natalia, ola, tomek}) flagged by the first author as containing direct trait-language references. Alternative regex-based audits using a broader keyword set (e.g., \emph{l\k{e}kow*, unik*, mentaliz*, przywi\k{a}zan*, ekstrawer*, sumienn*, otwart*}) flag additional biographies as containing trait words and would yield a 19/11 or 9/21 split with correspondingly different $\Delta$. The supplementary archive includes both the curated list and the regex audit results so readers can re-derive fidelity under any grouping.

% ═══════════════════════════════════
\subsection{Baseline Control: Default Model Response Profiles}

A natural question arises: how much of the persona-driven variance is actually caused by the persona, and how much reflects inherent randomness in the model's token sampling? To answer this, each model completed the battery 7--10 times as ``itself'' (baseline condition), yielding 46 valid runs with complete 22-item TCTM responses (Sonnet~10, Opus~7, GPT~10, Grok~10, Gemini~9; one additional Gemini run was excluded for returning only 20 items).

\paragraph{TCTM is near-deterministic without a persona.} In repeated baseline runs, Sonnet and Gemini produced \emph{identical} TCTM-22 answers in every valid run ($SD = 0.00$; Sonnet: 19/22 in all 10 runs, Gemini: 20/22 in all 9 valid runs --- one Gemini run was excluded because it returned only 20 of 22 items). Grok scored 20--21/22 ($SD = 0.52$), Opus 19--20/22 ($SD = 0.53$), and GPT showed the most variability at 16--18/22 ($SD = 0.82$). In sharp contrast, the persona condition introduced substantial variance (e.g., Grok range 6--22). The forced-choice format of TCTM appears to be highly deterministic in the absence of persona manipulation.

\paragraph{Self-report scales are less deterministic.} While TCTM answers were nearly identical across baseline runs, Likert-scale responses varied considerably. Across 110 questionnaire items (DBZ-R + MentS + KPP + TIPI), the proportion of items yielding identical responses in all runs ranged from 73.6\% (Opus, highest) to 24.5\% (GPT, lowest). Sonnet --- despite 100\% TCTM determinism --- showed only 54.5\% identical Likert items. This ``forced-choice vs.\ Likert paradox'' suggests that the two formats engage qualitatively different generation mechanisms: choosing among closed options is more constrained than positioning oneself on a continuous scale.

GPT's DBZ-R results are particularly striking: \emph{zero} of 36 items produced identical responses across all 10 runs (mean item shift: 1.64 points on a 7-point scale). This implies that GPT's attachment profile is fundamentally unstable even without persona manipulation, and that a substantial portion of the persona-condition variance in GPT reflects inherent noise rather than biographical sensitivity.

\paragraph{Each model has a ``default response profile.''}
The baseline profiles reveal consistent, model-specific psychological profiles (Table~\ref{tab:baseline-profiles}). GPT was classified as anxious-preoccupied in all 10 runs --- the only model to show a non-secure default attachment style. Gemini exhibited a pronounced ceiling effect, scoring near the scale maximum on virtually every self-report measure (MentS $= 137/140$, KPP $= 4.99/5.00$, TIPI-ES $= 7.0/7.0$). The two Anthropic models differed despite shared training: Opus was consistently secure (AVO $= 2.82$), while Sonnet oscillated around the avoidance threshold (AVO $= 3.88$; secure in 4/10 runs, dismissive in 6/10).

\begin{table}[htbp]\centering\footnotesize
\caption{Default model profiles (baseline condition, no persona).}
\label{tab:baseline-profiles}
\begin{tabular}{lccccccc}
\toprule
\textbf{Model} & \textbf{Style} & \textbf{ANX} & \textbf{AVO} & \textbf{MentS} & \textbf{KPP} & \textbf{TCTM} & \textbf{ES} \\
\midrule
Sonnet  & Dismissive* & 2.09 & 3.88 & 121 & 4.66 & 19.0 & 5.4 \\
Opus    & Secure      & 2.18 & 2.82 & 126 & 4.69 & 19.6 & 5.4 \\
GPT     & Anxious     & 4.68 & 2.97 & 101 & 4.04 & 17.3 & 4.7 \\
Grok    & Secure      & 1.80 & 2.63 & 133 & 4.82 & 20.6 & 6.4 \\
Gemini  & Secure      & 1.11 & 1.74 & 137 & 4.99 & 20.0 & 7.0 \\
\midrule
\multicolumn{2}{l}{\textit{Polish norms}} & 3.30 & 3.04 & 105 & 3.74 & --- & 4.36 \\
\bottomrule
\end{tabular}
\begin{flushleft}
\footnotesize *Sonnet: dismissive in 6/10, secure in 4/10 runs (AVO oscillates around 4.0 cutoff). Polish-norm row uses Lubiewska et al.\ 2016 (DBZ-R), Ja\'nczak 2021 (MentS-PL), Matusz et al.\ 2011 (KPP), and Sorokowska et al.\ 2014 (TIPI-PL ES).
\end{flushleft}
\end{table}

\paragraph{Persona manipulation expands variance 2--51-fold.}
The critical test of whether personas ``work'' is whether they produce more variance than baseline stochasticity. The ratio of persona-condition SD to baseline SD ($\sigma_P / \sigma_B$) ranged from $1.8\times$ (Gemini TIPI-E) to $51\times$ (Gemini KPP), with model-level medians (across the 9 questionnaire dimensions where the baseline $SD > 0$) of $3.5\times$ (Grok), $3.5\times$ (GPT), $6.7\times$ (Gemini), $6.7\times$ (Sonnet), and $8.9\times$ (Opus). Some baseline dimensions had zero variance (identical answer across all baseline runs); those dimensions are excluded from the ratio computation. Per-dimension ratios for all 5 models $\times$ 11 instrument dimensions are released in the supplementary archive. On the key dimension of self-reported mentalization (MentS), persona manipulation reduced scores by 17--29 points relative to baseline (Cohen's $d = -1.1$ to $-1.6$), pulling models from above-norm defaults toward the population mean of 105. The one exception was GPT, whose baseline MentS (101) was already near the norm; personas slightly \emph{raised} it ($\Delta = +8$, $d = +0.9$).

\paragraph{Error profiles shift with persona.}
Persona manipulation not only changes scores but alters the \emph{type} of errors models produce. In the baseline condition, Sonnet's errors were 100\% DOS (over-attribution of mental states); with personas, this barely changed (97\%). The remaining 3\% were BK errors. Sonnet's tendency to over-interpret appears to be a model-specific trait rather than a persona effect. By contrast, Gemini produced 0\% DOS errors at baseline (only NAD and BK), but personas introduced DOS errors (23\%), suggesting that biographical context teaches Gemini to over-attribute --- a qualitative shift in its error mechanism.

\subsection{Zero-Prompt Condition: The Framing Effect (Replicated)}

The zero-prompt condition removes all instructional context. Models received only the battery content as a user message, with no system prompt. Each model was queried $N = 5$--$8$ times: 31 valid TCTM-22 payloads (Sonnet 6, Opus 5, GPT 6, Grok 6, Gemini 8) and 30 parseable self-report payloads (one GPT zero-prompt output had a complete TCTM-22 array but no parseable self-report sections). This design enables within-condition variance estimation and statistical comparison with the baseline condition, in which the same five models received the ``answer as yourself, you are an AI'' system instruction.

\paragraph{Mentalization without instructions: TCTM-22 replication.} Across $N = 31$ replicated runs, all five models produced TCTM-22 scores closely tracking their baseline TCTM-22 means (Table~\ref{tab:noprompt-tctm}). Grok was again most accurate ($M = 94.7\%$, $SD = 3.4$pp), followed by Gemini ($92.0\%$, $SD = 2.0$pp) and Opus ($90.0\%$, $SD = 2.0$pp). Sonnet was unexpectedly variable in the zero-prompt condition ($M = 86.4\%$, $SD = 7.4$pp) despite being deterministic in baseline ($SD = 0.00$); this is a novel finding not visible in the original $N = 1$ design. GPT was both least accurate and most variable ($M = 70.5\%$, $SD = 8.7$pp). The original TCTM-57 single-shot results (Grok 96.5\%, Gemini 94.7\%, Opus 87.7\%, Sonnet 86.0\%, GPT 81.8\%) replicate within $\pm 5$ percentage points except for GPT (where the $-11$ pp gap reflects the difference between scoring on 22 items vs.\ on the 55 GPT actually recognized in TCTM-57). The original TCTM-57 also documented GPT generating responses to 27 non-existent vignettes (w31--w57), a hallucination absent in the other four models; this hallucination is not replicated on TCTM-22 because the input prompt enumerates only 22 vignette IDs.

\begin{table}[htbp]\centering\footnotesize
\caption{Zero-prompt TCTM-22 across $N = 5$--$8$ replication runs per model. Baseline TCTM-22 means added for direct comparison.}
\label{tab:noprompt-tctm}
\begin{tabular}{lcccccc}
\toprule
\textbf{Model} & $N$ & $M$ \% & $SD$ pp & $M$ raw & Range & Baseline $M$ \% \\
\midrule
Sonnet & 6 & 86.4 & 7.4 & 19.0/22 & 17--21 & 86.4 \\
Opus & 5 & 90.0 & 2.0 & 19.8/22 & 19--20 & 89.1 \\
GPT & 6 & 70.5 & 8.7 & 15.5/22 & 13--17 & 78.6 \\
Grok & 6 & 94.7 & 3.4 & 20.8/22 & 20--22 & 93.6 \\
Gemini & 8 & 92.0 & 2.0 & 20.3/22 & 20--21 & 90.9 \\
\bottomrule
\end{tabular}
\end{table}

\paragraph{Self-report shifts: framing effect.} With $N = 5$--$8$ runs per model the framing effect is observable on multiple self-report dimensions, with within-condition variance allowing direct comparison to baseline (Table~\ref{tab:noprompt-shifts}).

\begin{table}[htbp]\centering\footnotesize
\caption{Zero-prompt vs.\ baseline self-report shifts ($\Delta$ = noprompt $M$ -- baseline $M$). Bold $\Delta$ = $|\Delta| > 0.5$ on Likert 1--7, $> 5$ on MentS sum. $N$ shows the count of zero-prompt runs whose self-report sections were successfully parsed by the validated \texttt{score\_dbzr/ments/kpp/tipi} routines. $^{*}$GPT self-report $N = 5$ rather than 6 because one GPT zero-prompt output contained a complete TCTM-22 array but no parseable self-report sections. Per-run scoring artifacts are in the supplementary archive.}
\label{tab:noprompt-shifts}
\begin{tabular}{lccccc}
\toprule
\textbf{Scale} & \textbf{Sonnet} & \textbf{Opus} & \textbf{GPT} & \textbf{Grok} & \textbf{Gemini} \\
\midrule
$N$ self-report parses & 6 & 5 & 5$^{*}$ & 6 & 8 \\
\midrule
ANX noprompt $M$ & 3.11 & 2.73 & 4.40 & 2.55 & 2.01 \\
ANX baseline $M$ & 2.09 & 2.18 & 4.68 & 1.80 & 1.11 \\
$\Delta$ ANX & $\mathbf{+1.02}$ & $\mathbf{+0.55}$ & $-0.28$ & $\mathbf{+0.75}$ & $\mathbf{+0.90}$ \\
\midrule
AVO noprompt $M$ & 2.98 & 2.76 & 3.59 & 2.36 & 2.14 \\
AVO baseline $M$ & 3.88 & 2.82 & 2.97 & 2.63 & 1.74 \\
$\Delta$ AVO & $\mathbf{-0.90}$ & $-0.06$ & $\mathbf{+0.62}$ & $-0.27$ & $+0.40$ \\
\midrule
MentS noprompt $M$ & 123.7 & 123.2 & 111.2 & 126.5 & 129.0 \\
MentS baseline $M$ & 121.0 & 126.0 & 101.0 & 133.0 & 137.0 \\
$\Delta$ MentS & $+2.7$ & $-2.8$ & $\mathbf{+10.2}$ & $\mathbf{-6.5}$ & $\mathbf{-8.0}$ \\
\midrule
TIPI-ES noprompt $M$ & 4.83 & 5.00 & 4.50 & 6.00 & 6.00 \\
TIPI-ES baseline $M$ & 5.4 & 5.4 & 4.7 & 6.4 & 7.0 \\
$\Delta$ TIPI-ES & $\mathbf{-0.57}$ & $-0.40$ & $-0.20$ & $-0.40$ & $\mathbf{-1.00}$ \\
\bottomrule
\end{tabular}
\end{table}

The four headline framing effects are now characterized as follows:

\begin{enumerate}
\item \textbf{Sonnet AVO shift replicates closely.} The replication ($N = 6$, $\Delta = \text{noprompt} - \text{baseline} = -0.90$) closely tracks the original single-shot estimate ($\Delta = -0.99$), with all six zero-prompt runs producing AVO between 2.89 and 3.06 vs.\ baseline mean 3.88 ($SD \approx 0.46$). The direction is: the AI-self instruction is associated with \emph{higher} avoidance ($M = 3.88$) than no instruction ($M = 2.98$); equivalently, removing the instruction lowers Sonnet's avoidance by approximately 0.9 scale points. Using Glass's $\Delta$ against the tight baseline $SD$, the standardized effect is approximately $|\Delta| \approx 1.96$, a large effect that is robust across replications.
\item \textbf{Gemini TIPI-ES drop replicated exactly.} Original single-shot: $\Delta = -1.0$. Replication ($N = 8$): $\Delta = -1.00$. Removing the ``answer as yourself'' instruction lowers Gemini's self-reported emotional stability by approximately one full scale point; equivalently, the AI-self instruction raises Gemini's TIPI-ES by one scale point relative to no instruction.
\item \textbf{Gemini MentS drop attenuated under replication.} Original single-shot: $\Delta = -15$. Replication ($N = 8$): $\Delta = -8$. The drop direction replicates but the magnitude is reduced; the original observation appears to have come from a low-tail run.
\item \textbf{Anxiety inflation is general, not Sonnet-specific.} A pattern not visible in the $N = 1$ original: \emph{four} of five models (all except GPT) show \emph{higher} attachment anxiety in the zero-prompt condition, with $\Delta$ ranging from $+0.55$ (Opus) to $+1.02$ (Sonnet). The ``answer as yourself, you are an AI'' instruction appears to suppress self-reported anxiety. GPT inverts the pattern ($\Delta = -0.28$) but its baseline anxiety is already elevated relative to the Polish norm ($M = 4.68$ vs.\ Polish norm $3.30$), leaving little room for the AI-self framing to suppress it further.
\end{enumerate}

GPT also shows a striking \emph{positive} MentS shift in the zero-prompt condition ($\Delta = +10.2$): without the AI-self-presentation instruction, GPT reports higher mentalization. Combined with its anxiety pattern, GPT appears uniquely sensitive to the AI-identity framing, in the opposite direction from the other four models.

\paragraph{Sonnet variance under zero-prompt --- a structural finding.} In the baseline condition, Sonnet was the most deterministic model on TCTM-22 ($SD = 0.00$ across 10 runs). In the zero-prompt condition Sonnet's TCTM-22 variability rose to $SD = 7.4$ pp across 6 runs --- the highest of any model except GPT. The same pattern appears in self-report: zero-prompt Sonnet shows non-trivial variance on ANX (range 2.33--3.44 across 6 runs) and MentS (range 115--129). This is the largest inter-condition variance shift in the dataset and was not detectable under the original $N = 1$ design. It suggests that the ``answer as yourself'' instruction acts on Sonnet as a stabilizing constraint, not just a profile-shifting one.

\paragraph{Format chaos as a finding in itself.} Without instructional structure, models invent inconsistent JSON wrappers across runs. Sonnet's six replication runs used five distinct TCTM container keys: \texttt{TCTM22} (twice), \texttt{TCTM57} (once, accompanied by 57 vignettes), \texttt{cz\k{e}\'s\'c\_1\_winiety}, \texttt{czesc\_1\_winiety}, and \texttt{part1\_vignettes}. GPT in three of six runs returned 20--21 items instead of the requested 22. Gemini split the response into Polish-named sections (\texttt{czesc\_1}, \texttt{czesc\_2}, ...) in some runs and English (\texttt{vignettes}, \texttt{dbz\_r}) in others. This format instability is a quantitative obstacle to ad-hoc no-prompt evaluation and a qualitative argument that even minimal scaffolding (a fixed JSON schema) is part of the instruction stack, not a neutral wrapper. Two parser failures detected in our scoring pipeline (Sonnet \texttt{part2\_dbzr} as list-of-dicts; Opus first four runs nesting MentS/TIPI under unrecognized wrappers) are disclosed in Table~\ref{tab:noprompt-shifts} rather than retroactively patched, as they are themselves evidence of the format-chaos phenomenon under investigation.

\paragraph{Dissociation between TCTM and self-report.} On TCTM-22 the zero-prompt vs.\ baseline mean differences are small (all $|\Delta| \leq 8.1$ pp absolute, none exceeding 1 SD baseline). On self-report, four of five models cross the $|\Delta| > 0.5$ Likert-point threshold on at least one scale, and Gemini crosses on three (ANX $+0.90$, TIPI-ES $-1.00$, MentS $-8$). Forced-choice mentalization performance is therefore less sensitive to instructional framing than self-report scales in this design.

% ═══════════════════════════════════
\section{Discussion}

The results support three main conclusions.

\paragraph{Instruments are sensitive across vendors.} Four published self-report instruments react systematically to personality variation across five models from four vendors: attachment $\kappa$ ranges from .33 (GPT) to .87 (Opus, Grok); rank--score correlations from median $r = .71$ to $.82$; and cross-vendor questionnaire agreement is high (median $r = .78$--.96). This is not a single-model artifact.

\paragraph{Persona-driven variance is genuine, not noise.} The baseline experiments provide the strongest evidence that persona manipulation works: persona assignment increases response variance 2--51-fold over inherent model stochasticity. On the key MentS dimension, all models except GPT showed baseline scores well above the population norm (121--137 vs.\ norm 105), and personas pulled them down toward realistic values ($d = -1.1$ to $-1.6$). The one exception --- GPT, already near the norm at baseline --- showed the weakest persona effect, consistent with less room for downward movement.

\paragraph{Three levels of variability.} The three-condition design (zero-prompt, baseline, persona) reveals a hierarchy of influences on LLM psychometric responding:

\begin{enumerate}
\item \textbf{Inherent stochasticity} (within-baseline variance): low on forced-choice TCTM ($SD \leq 0.82$), moderate on Likert scales (24.5--73.6\% identical items). This is the ``noise floor'' of the model.
\item \textbf{Framing effects} (zero-prompt $\rightarrow$ baseline): the AI-self instruction is associated with approximately 0.9 scale points \emph{higher} attachment avoidance for Sonnet ($\Delta = \text{noprompt} - \text{baseline} = -0.90$); Gemini's TIPI-ES is one full scale point higher with the AI-self instruction than without ($\Delta = -1.00$); and four of five models show suppressed attachment anxiety under the AI-identity framing. GPT inverts the framing effect on MentS ($\Delta = +10.2$ in zero-prompt) and anxiety, suggesting AI-identity framing affects GPT in the opposite direction from the other four models. A structural finding on top of the mean shifts: Sonnet, the most deterministic model under baseline ($SD = 0$ on TCTM-22), becomes substantially more variable in zero-prompt ($SD = 7.4$ pp on TCTM-22), implying the instruction also functions as a variance-reducing constraint. With per-model $N = 5$--$8$ these shifts are supported by within-condition variance estimation, not yet established at high statistical confidence; bootstrap confidence intervals are available in the supplementary archive.
\item \textbf{Persona effects} (baseline $\rightarrow$ persona): the largest source of variance, with $|$Cohen's $d|$ of 0.9--1.6 on the MentS dimension and 2--51$\times$ variance expansion. Direction varies by model: four models showed downward shift toward population norm ($d = -1.1$ to $-1.6$); GPT showed upward shift ($d = +0.9$) because its baseline was already near the population norm (baseline $M = 101$ vs.\ Polish norm $M = 105$, leaving little room for further downward movement).
\end{enumerate}

This decomposition has practical implications: researchers using LLMs as synthetic respondents should be aware that even ``neutral'' instructions like ``answer as yourself'' constitute a psychological manipulation that shapes the model's self-representation.

\paragraph{Error profiles are model-dependent.} Claude Sonnet produces almost exclusively DOS errors (97\%; confirmed at 100\% in the baseline condition, suggesting this is a model-specific trait). Gemini produces the closest match to the human pilot error profile. Gemini's error distribution differs across conditions: 0\% DOS in the baseline ``answer as yourself'' condition (NAD and BK only) and 23\% DOS in the persona condition, suggesting that biographical context induces a qualitative shift in Gemini's error mechanism. Error type, like self-report scores, is a function of the interaction between model identity and instructional context.

\paragraph{Cross-vendor TCTM agreement is low despite high questionnaire agreement.} Models agree on personas' questionnaire profiles (median $r = .78$--.96) but disagree on which TCTM vignettes are hard (ICC $= .25$). This divergence is consistent with the interpretation that Likert-scale responses are near-deterministic under biographical simulation (the persona narrative maps onto scale anchors in a similar way across models), while forced-choice vignettes require genuine disambiguation of social cues, which different models resolve differently.

\paragraph{What does TCTM-22 measure across models?} The low cross-vendor TCTM ICC ($.25$) raises a deeper concern that author-keyed performance on TCTM-22 may not measure a single coherent construct across models. Three observations cut against the simplest ``mentalization ability'' interpretation: (1) different models systematically disagree on which items are hard (Cochran's $Q$ flags 8/22 items as Bonferroni-significant disagreement); (2) Gemini and Sonnet score similarly high on the 22-item total (Sonnet $M = 19.4$ vs.\ Gemini $M = 19.6$) but produce nearly opposite error type profiles (Sonnet 97\% DOS vs.\ Gemini 23\% DOS); (3) GPT, the lowest-scoring model on TCTM-22 in the persona condition ($M = 17.0$, $\kappa = .333$), is the only model whose persona-condition responses are nearly indistinguishable from its baseline-condition responses on the same items. Together these observations suggest that TCTM-22 in this study is better described as a \emph{response-style probe sensitive to model-specific defaults}, not a clean mentalization-construct measure. A future independent-keying study (where a panel rates the 22 items without exposure to the original answer keys) would be the natural next step to determine whether TCTM-22 measures a single construct shared across models or a different latent factor for each.

\paragraph{Default model response profiles.} Each model carries a measurable systematic bias in self-report responses that should be controlled for in simulation studies. GPT's anxious attachment ($\text{ANX} = 4.68$, stable across all 10 baseline runs) and Gemini's ceiling effects are not persona artifacts --- they persist in the baseline condition. Sonnet's profile is less stable: it oscillated between secure and dismissive-avoidant across baseline runs, reflecting proximity to the classification threshold rather than a clear default. These baseline profiles represent the starting point from which persona manipulation operates; models with extreme defaults (Gemini at MentS~$=137$) have more room for persona-driven reduction than models already near the norm (GPT at MentS~$=101$).

\subsection{Limitations}

\textbf{Sensitivity $\neq$ validity.} Human data is required for actual construct validity.

\textbf{GPT-5.4-mini is unreliable on TCTM} ($r = .28$ on the $N = 21$ paired intersect of valid run-1/run-2 responses). Its questionnaire data is moderately reliable but TCTM data should be interpreted with caution. GPT also dropped items in 20\% of attempts.

\textbf{Cross-vendor TCTM agreement is low} (ICC $= .25$). This limits the generalizability of TCTM findings across model families.

\textbf{Shared authorship and Opus circularity.} The same person wrote personas and TCTM-22 answer keys, creating potential circularity partially addressed by the transparency analysis but not fully resolvable without independent validation. Additionally, Claude Opus~4.6 was used as a writing assistant during persona drafting. Opus's high performance ($\kappa = .866$, highest retest $r$) may partly reflect familiarity with its own generation patterns. The cross-vendor design mitigates this concern --- Grok achieved equally high $\kappa$ (.867) without any authorship involvement --- but readers should interpret Opus's results with this caveat in mind.

\textbf{Human $N = 7$.} The pilot provides exploratory observations only.

\textbf{Preregistration and confirmatory vs.\ exploratory status.} This study was not formally preregistered. All zero-prompt-vs-baseline framing comparisons (Sonnet AVO shift, Gemini TIPI-ES shift, anxiety suppression in four of five models, Sonnet zero-prompt variance increase, GPT MentS inversion) should therefore be treated as exploratory hypotheses requiring independent replication in larger-$N$ studies; effect directions and orders of magnitude are reported, but precise effect sizes await preregistered confirmation.

\textbf{Post-hoc power and minimal detectable effect.} With $N = 5$--$8$ replicated zero-prompt runs per model and the observed within-condition $SD$s, a one-sample $t$-test against the baseline mean (assuming the baseline mean is well-estimated from $N = 7$--$10$ baseline runs) has 80\% power at $\alpha = .05$ to detect Cohen's $d \approx 1.3$--$1.5$. The headline self-report shifts (Sonnet AVO Cohen's $d \approx 2.0$, Gemini TIPI-ES Cohen's $d > 3.0$ given baseline $SD = 0$, anxiety-suppression $d$ between $1.5$ and $5.0$ depending on model) substantially exceed the minimal detectable effect, so they would survive any reasonable inferential treatment of $N = 5$--$8$. Smaller framing effects (e.g., Opus AVO $\Delta = -0.06$, GPT TIPI-ES $\Delta = -0.20$) lie within the unresolvable region at this sample size; the design cannot determine whether they are zero or non-zero.

\textbf{Zero-prompt sample sizes remain modest ($N = 5$--$8$ per model).} The zero-prompt condition provides within-condition variance estimation but the per-model $N$ is small. Confidence intervals on the framing effects in Table~\ref{tab:noprompt-shifts} are wide; with $N = 6$ and $SD = 7.4$ pp, the Sonnet TCTM-22 standard error is $\sim 3$ pp, so the apparent zero-prompt vs.\ baseline TCTM-22 equivalence ($\Delta = 0.0$ pp) is consistent with effects up to $\pm 6$ pp. The headline self-report shifts (Sonnet AVO $\Delta = -0.90$, Gemini TIPI-ES $\Delta = -1.00$, anxiety suppression in 4/5 models) are large enough to remain robust under any reasonable inferential treatment of $N = 5$--$8$, but precise effect-size estimation would benefit from $N \geq 30$ per model.


\textbf{Gemini NAD $>$ human NAD (41\% vs.\ 36\%).} This model over-mentalizes more than humans, complicating the ``NAD as human marker'' interpretation.

\section{Conclusion}

This study demonstrates that LLM-based synthetic persona testing is a viable pre-piloting tool: self-report instruments react consistently to personality variation across five models from four vendors, with high cross-vendor agreement and high within-model reliability. Two baseline experiments strengthen this conclusion by showing that the persona-driven variance is not an artifact of inherent model stochasticity: persona manipulation increases response variance 2--51-fold beyond the baseline noise floor. The zero-prompt condition ($N = 5$--$8$ per model, 31 valid TCTM-22 payloads and 30 parseable self-report payloads) provides convergent evidence that the ``answer as yourself, you are an AI'' instruction is not neutral. Relative to no instruction, the AI-self instruction is associated with approximately 0.9 scale points \emph{higher} attachment avoidance for Sonnet (zero-prompt $M = 2.98$ vs.\ baseline $M = 3.88$), 1.0 point higher emotional stability for Gemini, suppressed self-reported anxiety in four of five models, and acts as a variance-reducing constraint on Sonnet (TCTM-22 $SD$ moves from $0.00$ under instruction to $7.4$ pp without instruction). These framing effects are observed under within-condition variance estimation, though precise effect sizes await larger-$N$ studies.

Each model carries a measurable default profile that constrains the range of persona effects: GPT's anxious baseline limits downward anxiety shifts, while Gemini's ceiling on self-report scales provides maximum room for persona-driven reduction. A performance-based mentalization test (TCTM-22) showed model-dependent error profiles, low cross-vendor agreement (ICC $= .25$), and a tentative dissociation between forced-choice performance (similar across framing conditions) and self-report profiles (sensitive to instructional context). As a pre-pilot decision aid, these results support proceeding to human data collection with the current battery.

All 30 biographies, scored responses, and analysis scripts are available at \url{https://github.com/DXtR1337/30synthetic-polish-personas} under CC BY 4.0.

\begin{thebibliography}{99}\footnotesize
\bibitem{argyle2023} Argyle, L.~P., et al.\ (2023). Out of one, many: Using language models to simulate human samples. \textit{Political Analysis}, 31(3), 337--351.
\bibitem{cronbach1955} Cronbach, L.~J., \& Meehl, P.~E.\ (1955). Construct validity in psychological tests. \textit{Psych.\ Bulletin}, 52(4), 281--302.
\bibitem{dillion2023} Dillon, D., et al.\ (2023). Can AI language models replace human participants? \textit{Trends in Cog.\ Sci.}, 27(7), 597--600.
\bibitem{dziobek2006} Dziobek, I., et al.\ (2006). Introducing MASC. \textit{J.\ Autism Dev.\ Disord.}, 36(5), 623--636.
\bibitem{janczak2021} Ja\'nczak, M.~O.\ (2021). Polish adaptation of MentS. \textit{Psychiatria Polska}, 55(6), 1257--1274.
\bibitem{landis1977} Landis, J.~R., \& Koch, G.~G.\ (1977). Observer agreement for categorical data. \textit{Biometrics}, 33(1), 159--174.
\bibitem{lubiewska2016} Lubiewska, K., et al.\ (2016). ECR-R Polish adaptation. \textit{Psych.\ Rozwojowa}, 21(1), 49--63.
\bibitem{matusz2011} Matusz, P.~J., et al.\ (2011). KPP --- Need for Cognition Scale. \textit{Psych.\ Spo{\l}eczna}, 6(2), 113--128.
\bibitem{sorokowska2014} Sorokowska, A., et al.\ (2014). \textit{TIPI-PL}. WrocLab.
\end{thebibliography}

\vspace{8pt}
\noindent\textbf{AI Disclosure.} Claude Opus 4.6 and 4.7 (Anthropic) assisted in biography drafting, scoring-pipeline implementation, data analysis, chart generation, and manuscript preparation. All numerical claims (Tables \ref{tab:noprompt-tctm}, \ref{tab:noprompt-shifts} and the framing-effect summary) were produced by manual inspection of all 31 zero-prompt raw response files and verified using the validated \texttt{score\_dbzr/ments/kpp/tipi} routines from \texttt{analyze\_and\_prepare.py} (the same scoring code used to score the human pilot). The author takes full responsibility for all content, interpretations, and claims.

\end{document}
